\begingroup
\makeatletter\tagsleft@true\let\veqno\@@leqno\makeatother
\clearpage\NoSourceHead\SourceLeaf{339}
\SourceGroup{II}
\begin{center}\large\bfseries [ALLGEMEINES ÜBER DIE REIHE $F(\alpha,\beta,\gamma,x)$.]\end{center}
\HistoricCitation{[Aus Handbuch 18, Bd, Mathematische Brouillons, October 1805.]}
\NumberedItem{1}
\noindent{\small[S. 234]}\par
\par\indent Es sei
\[1+\frac{\alpha{.}\beta}{1{.}n}x+\frac{\alpha{.}\alpha+1{.}\beta{.}\beta+1}{1{.}2{.}n{.}n+1}xx+{.}{.}{.}=\{\alpha,\beta,n\}=P\]
\begin{equation}
\{\alpha,\beta,n+1\}=\{\alpha,\beta,n\}-\frac{\alpha\beta x}{n{.}n+1}\{\alpha+1,\beta+1,n+2\},\tag*{1)}
\end{equation}
\begin{equation}
\{\alpha,\beta+1,n\}=\{\alpha,\beta,n\}+\frac{\alpha x}{n}\{\alpha+1,\beta+1,n+1\},\tag*{2)}
\end{equation}
\begin{equation}
\{\alpha,\beta+1,n+1\}=\{\alpha,\beta,n\}+\frac{\alpha(n-\beta)x}{n{.}n+1}\{\alpha+1,\beta+1,n+2\}.\tag*{3)}
\end{equation}
\noindent Also
\begin{equation}
(n-\beta)\{\alpha,\beta,n+1\}+\beta\{\alpha,\beta+1,n+1\}=n\{\alpha,\beta,n\},\tag*{4)}
\end{equation}
\begin{equation}
\{\alpha,\beta+1,n\}=\{\alpha+1,\beta,n\}+\frac{(\alpha-\beta)x}{n}\{\alpha+1,\beta+1,n+1\},\tag*{5)}
\end{equation}
\begin{equation}
\beta\{\alpha,\beta+1,n\}=\alpha\{\alpha+1,\beta,n\}-(\alpha-\beta)\{\alpha,\beta,n\}.\quad\text{gut.}\tag*{6)}
\end{equation}
\NumberedItem{2}
\par\indent Wenn in der D[ifferential] G[leichung]
\[\begin{aligned}
&(x-xx)d^2P-(x-xx)\frac{dP\,d^2x}{dx}\\[.6em]
&\quad+\{\gamma-(\alpha+\beta+1)x\}\,dPdx-\alpha\beta Pdx^2\ [=0]
\end{aligned}\]
\noindent statt $[x]$
\clearpage\SourceLeaf{340}\ArithmeticHead{333}{R}{{\scriptsize ZUR THEORIE DER REIHE $F(\alpha,\beta,\gamma,x)$.\quad II.}}
\[x=\frac{t}{t-1}\]
\noindent subst[ituirt] wird, so erhält man[*)
\[\begin{aligned}
&(1-t)(t-t^2)ddP+(1-t)(\gamma+(\alpha+\beta-1-\gamma)t)\,dPdt\\[.4em]
&\quad+\alpha\beta Pdt^2=0.
\end{aligned}\]
\noindent Ferner ergibt sich für]
\[\begin{aligned}
x&=\frac{4t}{(1+t)^2},& dx&=\frac{4(1-t)dt}{(1+t)^3},\\[.7em]
ddx&=[-]\frac{8(2-t)dt^2}{(1+t)^4},& x-xx&=\frac{(1-t)^2[4t]}{(1+t)^4},
\end{aligned}\]
\[\begin{aligned}
&\frac{[4t](1-t)}{(1+t)^4}P''+\frac{[4t](4-2t)}{(1+t)^5}P'\\[.6em]
&\quad+\left\{\gamma-(\alpha+\beta+1)\frac{4t}{(1+t)^2}\right\}P'\frac4{(1+t)^3}
-\alpha\beta P\frac{16(1-t)}{(1+t)^6}\ [=0.]
\end{aligned}\]
\clearpage\NumberedItem{3}
\noindent{\small[S. 235]}\par
\[\frac{1+\frac{A{.}B+\beta}{1{.}C+\gamma}x+\frac{A{.}A+\alpha{.}B+\beta{.}B+2\beta}{1{.}2{.}C+\gamma{.}C+2\gamma}xx{.}{.}{.}}{1+\frac{A{.}B}{1{.}C}x+\frac{A{.}A+\alpha{.}B{.}B+\beta}{1{.}2{.}C{.}C+\gamma}xx{.}{.}{.}}\]
\[\begin{aligned}
&=\left\{\frac A\alpha,\ \frac B\beta,\ \frac C\gamma,\ \frac{\alpha\beta x}\gamma\right\}\ [=]\\[.8em]
&\quad\cfrac1{1-\cfrac{\frac{A{.}\beta C-\gamma B}{C{.}C+\gamma}x}{1-\cfrac{\frac{B+\beta{.}\alpha C+\alpha\gamma-A\gamma}{C+\gamma{.}C+2\gamma}x}{1-\cfrac{\frac{A+\alpha{.}\beta C+\beta\gamma-\gamma B}{C{.}C+\gamma}x}{\text{etc.}}}}}
\end{aligned}\]
\[C=1,\quad\gamma=0\]
\[\begin{aligned}
&\cfrac1{1-\cfrac{A\beta x}{1-\cfrac{(B+\beta)\alpha x}{1-\cfrac{(A+\alpha)\beta x}{1-\cfrac{(B+2\beta)\alpha x}{\begin{array}{@{}l@{}}{.}\\[-.5em]\kern.3em{.}\\[-.5em]\kern.6em{.}\\[-.5em]\kern.9em{.}\end{array}}}}}}\\[.8em]
&\quad=\frac{1+A(B+\beta)x+{.}{.}{.}}{1+ABx+{.}\quad{.}}
\end{aligned}\]
\[\operatorname{tang}[\varphi]=\cfrac\varphi{1-\cfrac{\frac{\varphi\varphi}3}{1-\cfrac{\frac{\varphi\varphi}{15}}{\text{etc.}}}}\]
\par\Needspace{5\baselineskip}\bigskip\noindent\rule{.25\linewidth}{.4pt}\par\smallskip
{\small\par\indent [*) In der Handschrift ist diese Transformation nicht ganz durchgeführt; das Ergebnis ist deshalb nach Werke III, S. 218 ergänzt.]\par}
\clearpage
\endgroup
